\documentclass{article}

\usepackage[french]{babel}
\usepackage[utf8]{inputenc}
\usepackage[T1]{fontenc}
\usepackage{amsmath,amssymb,amsthm}
\usepackage{graphicx}
\usepackage{enumitem}
\usepackage{hyperref}
\usepackage{bookmark}
\usepackage{geometry}
\usepackage{tikz}
\usetikzlibrary{positioning,chains,fit,shapes,calc}
\usepackage[most]{tcolorbox}
\usepackage[dvipsnames]{xcolor}
\geometry{
  a4paper,
  left=25mm,
  right=25mm,
  top=8mm,
  bottom=25mm}
\usepackage{fancyhdr}

% Suppression des alinéas automatiques
\setlength{\parindent}{0pt}

\begin{document}

\textit{(n°candidat) COADIC Alban}

\section*{Simulation d'écoulement de fluides sur des surfaces solides.}
\vspace{5mm}


L'écoulement d'un fluide autour d'un obstacle peut générer des structures
tourbillonnantes appelées vortex de von Kármán. Ils sont à l'origine de structures
périodiques de tourbillons alternés induisant des forces cycliques qui représentent un
défi majeur pour l'intégrité des structures, pouvant engendrer des phénomènes de
résonance destructeurs.

\vspace{8mm}

Bien que prévisibles en théorie, les vortex de von Kármán créent des instabilités de
pression complexes derrière les obstacles. Leur maîtrise est cruciale pour l'optimisation
de structures soumises à des vents violents ou des courants marins. La conception d'une
modélisation en temps réel aiderait à la sécurisation des infrastructures



\section*{Positionnement thématique (ÉTAPE 1) :}
- \textit{PHYSIQUE} (Mécanique des fluides)\\
- \textit{INFORMATIQUE} (Modélisation et simulation numérique (CFD))\\
- \textit{PHYSIQUE} (Mécanique du solide)\\


\section*{Mots-clés (ÉTAPE 1) :}
\begin{tabular}{ll}
\textbf{Mots-Clés (français)} & \textbf{Mots-Clés (anglais)} \\
\textit{Vortex de Von Kármán} & \textit{Von Kármán vortex} \\
\textit{Equations de Navier-Stokes} & \textit{Navier-Stokes equations} \\
\textit{Simulation en temps réel} & \textit{Real-time simulation} \\
\textit{Résistance structurelle} & \textit{Structural resistance} 
\end{tabular}
\vspace{5mm}  


\section*{Bibliographie commentée}

La mise en place d'une simulation numérique efficace nécessite une solide maîtrise des 
fondements théoriques de la mécanique des fluides. Une formation universitaire spécialisée 
fournit les clés pour intégrer ces modélisations dans le contexte général de la dynamique 
des écoulements incompressibles et de la théorie de la couche limite [1]. Cette assise 
théorique est essentielle pour analyser les résultats obtenus par le simulateur et pour 
concevoir des géométries structurelles résistantes aux instabilités périodiques.\\

La simulation des fluides autour d'obstacles repose sur une mise en œuvre algorithmique 
efficace, particulièrement lorsqu'on vise une interaction en temps réel. Un article fondateur 
dans ce domaine propose une approche privilégiant la stabilité visuelle et la rapidité de 
calcul, permettant de simuler des écoulements complexes comme la fumée ou l'air de manière 
interactive [2]. Ce travail se concentre sur une résolution simplifiée des équations de 
Navier-Stokes [3], garantissant la conservation de la masse et de la quantité de mouvement.\\

Toutefois, pour une application rigoureuse en ingénierie, cette rapidité doit être confrontée 
à un choix. Il existe deux approches principales pour simuler un fluide. La méthode Eulérienne, qui 
considère le fluide à travers un maillage fixe et la méthode Lagrangienne, qui suit des particules
individuelles [4]. Le choix de la méthode influence directement la précision et la complexité
de la simulation, et afin d'être capable d'observer les structures périodiques formées il est
préférable d'opter pour une approche eulérienne[2].\\

Une fois la méthode choisie, la mise en œuvre numérique nécessite de traiter chaque 
composante de la grille indépendamment afin de satisfaire les équations de Navier-Stokes[3] à
tout instant. Il faut alors implémenter une méthode itérative parcourant chaque cellule de la
grille pour mettre à jour les valeurs de vitesse et de pression en fonction des conditions aux
limites et des interactions fluide-structure[2].\\

Par ailleurs, la conservation des grandeurs physiques, comme la masse ou le volume, reste un 
défi majeur dans les simulations à temps de calcul réduit. L'étude d'un schéma semi-lagrangien 
spécifique, le CIP-CSL, apporte une solution technique permettant de garantir une conservation 
exacte tout en maintenant une interface nette autour de l'obstacle [5]. Cette rigueur 
numérique est indispensable pour que les forces prédites par la simulation correspondent à la 
réalité physique des vortex de von Kármán.\\

Pour garantir la robustesse et la performance d'un tel outil, le choix du langage de 
programmation est déterminant. Il faut être en mesure de paralléliser efficacement les calculs. 
La documentation du langage Rust souligne comment la gestion stricte de la mémoire et la performance 
brute de ce langage permettent de développer des outils de calcul intensif fiables [6]. L'utilisation 
de Rust offre ainsi une alternative moderne aux langages classiques, assurant la stabilité du 
simulateur lors de ces calculs parallélisés.\\

  
\section*{Problématique retenue}


Comment est-il possible de modéliser et simuler en temps réel l'écoulement de fluides autour 
d'obstacles solides   afin de prédire les forces exercées et d'optimiser la conception des structures pour minimiser
les effets des vents et des vortex de von Kármán sur ces dernières ?\\



\section*{Objectifs du TIPE du candidat}

Simuler en temps réel l'écoulement de fluides autour d'obstacles solides.\\

Prédire les forces exercées par les vents sur les structures.\\

Visualiser les vortex de von Kármán générés par l'écoulement autour des obstacles.\\

Analyser l'impact des vortex et des vents sur la résistance structurelle et proposer des solutions pour minimiser ces effets.\\





\section*{Références bibliographiques (ÉTAPE 1)}

\textbf{[1]} Buffat, M. Cours de mécanique des fluides approfondie, 2023. \url{https://perso.univ-lyon1.fr/marc.buffat/COURS/BOOK_MECAFLU_HTML/intro.html}\\

\textbf{[2]} Stam, J. Real-Time Fluid Dynamics for Games, 2003. \textit{Proceedings of the Game Developer Conference}, 2003.\\

\textbf{[3]} Couderc, F. Développement d'un code de calcul pour la simulation d'écoulements de fluides non miscibles. Thèse de doctorat, Université Paul Sabatier Toulouse III, 2007.\\

\textbf{[4]} Resolved Analytics. Lagrangian vs Eulerian Phase Models, 2020. \url{https://www.resolvedanalytics.com/cfd-physics-models/lagrangian-vs-eulerian-phase-models}\\

\textbf{[5]} Yabe, T., Tanaka, R., Nakamura, T., Xiao, E. An Exactly Conservative Semi-Lagrangian Scheme (CIP-CSL) in One Dimension. \textit{Journal of Computational Physics}, 169, (2001) 556--593.\\

\textbf{[6]} Klabnik, S., Nichols, C. Le langage de programmation Rust, 2021. \url{https://jimskapt.github.io/rust-book-fr/title-page.html}\\






\end{document}